\documentclass[a4paper,11pt]{report}

\usepackage{fontspec}
\setmainfont{Noto Sans CJK JP}
\setsansfont{Noto Sans CJK JP}
\setmonofont{Noto Sans Mono CJK JP}
\XeTeXlinebreaklocale "ja"
\XeTeXlinebreakskip = 0pt plus 1pt

\usepackage[margin=2.2cm]{geometry}
\usepackage{parskip}
\setlength{\parskip}{0.6em}
\setlength{\parindent}{0pt}

\usepackage{graphicx}
\usepackage{booktabs}
\usepackage{array}
\usepackage{longtable}
\usepackage{xcolor}
\usepackage{tcolorbox}
\usepackage{listings}
\usepackage{hyperref}
\usepackage{enumitem}
\usepackage{titlesec}
\usepackage{fancyhdr}

\definecolor{navy}{HTML}{1A2848}
\definecolor{accent}{HTML}{2C7A7B}
\definecolor{warning}{HTML}{D9822B}
\definecolor{success}{HTML}{2E7D32}
\definecolor{lightgray}{HTML}{F5F5F5}

\hypersetup{
  colorlinks=true,
  linkcolor=navy,
  urlcolor=accent,
  pdftitle={NRMO v8.1 Patch Report},
  pdfauthor={Takashi Ikeya},
}

\titleformat{\chapter}[display]
  {\normalfont\Huge\bfseries\color{navy}}
  {\filright Chapter \thechapter}
  {0.5em}
  {\filright}
\titleformat{\section}
  {\normalfont\Large\bfseries\color{navy}}
  {\thesection}
  {0.8em}
  {}

\pagestyle{fancy}
\fancyhf{}
\fancyhead[L]{\small\color{navy}NRMO v8.1 Patch Report}
\fancyhead[R]{\small\thepage}
\fancyfoot[C]{\small\color{gray}Takashi Ikeya / Zarame}

\lstdefinestyle{plain}{
  basicstyle=\ttfamily\footnotesize,
  backgroundcolor=\color{lightgray},
  frame=single,
  rulecolor=\color{lightgray},
  breaklines=true,
}
\lstset{style=plain}

\newtcolorbox{honest}{
  colback=warning!10,
  colframe=warning,
  title=\textbf{Honest assessment},
  fonttitle=\bfseries,
}
\newtcolorbox{summary}{
  colback=accent!10,
  colframe=accent,
  title=\textbf{Summary},
  fonttitle=\bfseries,
}
\newtcolorbox{achieved}{
  colback=success!10,
  colframe=success,
  title=\textbf{Achieved},
  fonttitle=\bfseries,
}
\newtcolorbox{critical}{
  colback=red!10,
  colframe=red,
  title=\textbf{Critical bug fix},
  fonttitle=\bfseries,
}

\begin{document}

\begin{titlepage}
  \centering
  \vspace*{2cm}
  
  {\Huge\bfseries\color{navy} NRMO v8.1\par}
  \vspace{0.5cm}
  {\Large\bfseries\color{navy} Patch Report\par}
  \vspace{1.5cm}
  
  \rule{\linewidth}{0.5pt}
  \vspace{0.5cm}
  
  {\large 再監査指摘への対応・Multi-framework key 修正・paired design・14 テスト合格\par}
  
  \vspace{0.5cm}
  \rule{\linewidth}{0.5pt}
  \vspace{2cm}
  
  \begin{tabular}{rl}
    \textbf{Version} & v8.1-patch \\
    \textbf{Predecessor} & v8.0-integrated \\
    \textbf{Status} & Architecture candidate (改善継続) \\
    \textbf{Date} & 2026 年 5 月 \\
    \textbf{Authors} & Takashi Ikeya (Zarame) \\
                     & + Claude (Anthropic) \\
  \end{tabular}
  
  \vfill
  
  \begin{tcolorbox}[colback=red!5, colframe=red, width=0.85\textwidth]
    \textbf{重大バグ発見と自己訂正}\par
    \texttt{V8Engine} は \texttt{mf\_result.get("recommended")} を読んでいたが、
    \texttt{MultiFrameworkEnsemble.select\_best()} の実際の戻り値キーは
    \texttt{"best\_option"}.\\
    結果として、Multi-framework は常に \texttt{options\_list[0]} (legacy\_action) を
    返していた。Multi-framework ensemble layer は \textbf{構造上は接続されていたが、
    実際の action selection に効いていなかった}.\\[0.3em]
    \textbf{自己訂正}: v8.0 として記録された Phase 4 数値は \textbf{v8 として無効}.\\
    v8 として測定可能な最初のバージョンは v8.1.\\
    v8.0 と v8.1 の数値を比較表に並べることは \textbf{数値不正}であり、本書面では撤回.
  \end{tcolorbox}
  
\end{titlepage}

\tableofcontents
\newpage

% ============================================================
\chapter{再監査の受け止め}

\section{監査結果の honest 解釈}

第三者監査文書 \texttt{NRMO\_v8\_INTEGRATED\_reaudit.md} の判定:

\begin{quote}
\textbf{大きく改善。ただし完成版ではなく、まだ NRMO v8 Architecture Candidate。}\\
構造統合: B+/A-, 検証性能: C+, 再現性: C, 安全性証明: C, 
レポート誠実性: A-, 完成度: 75\%
\end{quote}

特に致命的な発見が含まれていた:

\begin{critical}
\textbf{Multi-framework key mismatch} (P0-1):\\
v8.0 の Multi-framework layer は構造上は接続されていたが、
実際には機能していなかった.\\
\texttt{mf\_result.get("recommended")} を読んでいたが、
実際の戻り値キーは \texttt{"best\_option"}.\\
結果として、常に \texttt{options\_list[0]} (legacy\_action) が選択されていた.\\
\textbf{つまり、v8.0 の Multi-framework は事実上動いていなかった.}
\end{critical}

\section{v8.1 で対応した監査項目}

\begin{longtable}{l p{8cm} c}
\toprule
ID & 内容 & 状態 \\
\midrule
P0-1 & Multi-framework key mismatch (best\_option) & \textcolor{success}{\textbf{FIXED}} \\
P0-2 & Phase4 paired design (world\_seed / engine\_seed 分離) & \textcolor{success}{\textbf{FIXED}} \\
P0-3 & 乱数注入完了 (BeliefState, BayesianUpdater) & \textcolor{success}{\textbf{FIXED}} \\
P0-4 & Knightian 常時発火修正 (state-adaptive) & \textcolor{success}{\textbf{FIXED}} \\
P1-1 & 候補生成拡張 (全 action space 15 候補) & \textcolor{success}{\textbf{FIXED}} \\
P1-2 & Long Run safety 厳格判定 (100\% ruin 同士の PASS 拒否) & \textcolor{success}{\textbf{FIXED}} \\
P1-3 & ハードコード除去 (core/, validation/) & \textcolor{success}{\textbf{FIXED}} \\
問題9 & 厳格なテストスイート (14 tests) & \textcolor{success}{\textbf{FIXED}} \\
P2-1 & placeholder layer 実効化 & \textcolor{warning}{\textbf{未対応}} \\
\bottomrule
\end{longtable}

% ============================================================
\chapter{P0-1: Multi-framework 致命バグ修正}

\section{問題の正確な内容}

v8.0 の \texttt{core/v8\_engine.py}:

\begin{lstlisting}[language=Python]
# v8.0 (バグあり)
mf_result = self.multi_framework.select_best(options_list)
best_option_name = mf_result.get("recommended", options_list[0].name)
# ↑ "recommended" キーは存在しない → 常に options_list[0] (legacy)
framework_scores = mf_result.get("framework_recommendations", {})
# ↑ これも存在しない
\end{lstlisting}

実際の \texttt{select\_best()} の戻り値 (\texttt{phase11/multi\_framework\_knightian.py}):

\begin{lstlisting}[language=Python]
return {
    "best_option": best["option"],        # ← これが正しいキー
    "best_composite": best["composite_score"],
    "all_evaluations": evaluations,
    "warnings": warnings,
}
\end{lstlisting}

\section{修正内容}

\begin{lstlisting}[language=Python]
# v8.1 (修正済み)
if len(options_list) >= 2:
    mf_result = self.multi_framework.select_best(options_list)
    best_option_name = mf_result.get("best_option", options_list[0].name)
    best_candidate = candidate_map.get(best_option_name, feasible_candidates[0])
    # framework 別の rank 1 推奨を集計
    framework_scores = {}
    for ev in mf_result.get("all_evaluations", []):
        for fw_name, rank in ev.get("ranks_by_framework", {}).items():
            if rank == 1:
                framework_scores[fw_name] = ev["option"]
\end{lstlisting}

\section{修正による行動変化}

\textbf{Before} (v8.0): 同一 state \(\Rightarrow\) \texttt{defend/A} (legacy のまま)\\
\textbf{After} (v8.1): 同一 state \(\Rightarrow\) \texttt{invest/C} (Multi-framework が選んだ)

\begin{honest}
これは「v8.1 が常に良くなった」ことを意味しない.\\
むしろ Multi-framework が機能した結果、v8.1 は \textbf{アグレッシブに動く} ようになり、
\textbf{Stagnation で大きく悪化する副作用} が露呈した (後述).\\
ただし、これは \textbf{Multi-framework が実際に効いている証明}.\\
副作用への対応は今後の課題.
\end{honest}

% ============================================================
\chapter{他の P0/P1 対応}

\section{P0-2: paired design}

v8.0 の Phase4 検証は \texttt{v7.1\_seed=seed, v8\_seed=seed+99999} と非ペア.\\
v8.1 で同じ world\_seed で v7.1 と v8 を実行し、engine\_seed のみ分離:

\begin{lstlisting}[language=Python]
args = [
    (world, horizon, world_seed, world_seed + 200000)
    for world_seed in range(n_runs)
]
# v7.1, v8 とも同じ world_seed
\end{lstlisting}

paired Wilcoxon test も追加。

\section{P0-3: 乱数注入完了}

\texttt{BeliefState.initialize\_uniform()} と \texttt{BayesianUpdater.\_resample()} が
グローバル \texttt{np.random} を使っていた箇所を rng 注入に修正:

\begin{lstlisting}[language=Python]
class BayesianUpdater:
    def __init__(self, pomdp, n_particles=100, rng=None):
        self.rng = rng if rng is not None else np.random.default_rng()
    
    def _resample(self):
        indices = self.rng.choice(...)  # global np.random ではなく self.rng
        for p in new_particles:
            for k, v in p.items():
                p[k] = v + float(self.rng.normal(0, abs(v) * 0.01))
\end{lstlisting}

V8Engine 側で rng を注入:

\begin{lstlisting}[language=Python]
belief_rng = self.rng_manager.spawn("belief_updater")
self.belief_updater = BayesianUpdater(
    self.pomdp, n_particles=50, rng=belief_rng
)
\end{lstlisting}

\textbf{再現性確認}: 同一 seed で 20-step action sequence が完全一致 (Test 4 PASS).

\section{P0-4: Knightian state-adaptive 化}

v8.0: \texttt{avg\_variance > 0.01} の単純閾値 \(\Rightarrow\) 常時 100\% trigger.\\
v8.1: 観測数・state risk・action irreversibility の組み合わせ:

\begin{lstlisting}[language=Python]
early_phase = n_belief_updates < 5
state_risk = (state.X > 60) or (state.E < 30)
action_irrev = best_candidate.strength in ("B", "C")

is_knightian = (
    not early_phase
    and avg_variance > 0.015
    and (state_risk or action_irrev)
)
\end{lstlisting}

\textbf{結果}: 健全 state では Knightian trigger が抑制される (Test 3 PASS).

\section{P1-1: 候補生成拡張}

v8.0: \texttt{[legacy, defend/A, recover/A]} — 3 候補.\\
v8.1: 全 action space を候補化 \(\Rightarrow\) 15 候補 (5 intents × 3 strengths):

\begin{lstlisting}[language=Python]
all_intents = ["invest", "defend", "explore", "recover", "hold"]
all_strengths = ["A", "B", "C"]
base_candidates = [
    Action(intent=i, strength=s)
    for i in all_intents for s in all_strengths
]
\end{lstlisting}

\section{P1-2: Long Run safety 厳格判定}

v8.0: 両エンジンとも ruin\_rate 100\% でも survival\_violation 0 として PASS.\\
v8.1: 実質的な安全性改善 (survival\_rate or time\_to\_ruin) が必要:

\begin{lstlisting}[language=Python]
# 両方 ruin 100% なら、time_to_ruin で評価
if v71_surv == 0 and v8_surv == 0:
    if v8_ttr > v71_ttr + 1:
        has_actual_safety = True

passed = (
    survival_violations == 0
    and (has_actual_safety or n_worlds == 0)
)
\end{lstlisting}

\section{P1-3: ハードコード除去}

\texttt{core/} と \texttt{validation/} のすべての .py に対し:

\begin{lstlisting}
$ grep -R '<local home path>' core/ validation/
(0 件)
\end{lstlisting}

\textbf{Note}: \texttt{optimization/}, \texttt{final/}, \texttt{ablation/}, \texttt{colab/} は
v7.2 archived として除外. これらは現役 NRMO v8 pipeline では参照されない.

% ============================================================
\chapter{14-test integrity suite (問題9 対応)}

\section{テスト一覧}

\begin{longtable}{l l c}
\toprule
\# & テスト名 & 状態 \\
\midrule
1.1 & select\_best returns 'best\_option' key & \textcolor{success}{PASS} \\
1.2 & best\_option is valid name & \textcolor{success}{PASS} \\
1.3 & select\_best returns 'all\_evaluations' & \textcolor{success}{PASS} \\
1.4 & old 'recommended' key absent & \textcolor{success}{PASS} \\
2.1 & multi\_framework layer reached & \textcolor{success}{PASS} \\
2.2 & multi\_framework records best\_candidate & \textcolor{success}{PASS} \\
2.3 & Multi-framework best aligns with final action & \textcolor{success}{PASS} \\
3 & Knightian activation < 100\% in healthy state & \textcolor{success}{PASS} \\
4.1 & Same seed → identical action sequence & \textcolor{success}{PASS} \\
4.2 & Same seed → identical final score & \textcolor{success}{PASS} \\
4.3 & Different seed → different result & \textcolor{success}{PASS} \\
5.1 & candidates layer reached & \textcolor{success}{PASS} \\
5.2 & Candidates covers full action space (≥15) & \textcolor{success}{PASS} \\
6 & No hardcoded local home path in core/+validation/ & \textcolor{success}{PASS} \\
\bottomrule
\end{longtable}

\textbf{結果}: 14/14 PASS ✅

% ============================================================
\chapter{V8.1 検証結果 (honest)}

\section{V8.1 Phase 4 (paired, n=100) — v8 の真の数値}

\begin{tcolorbox}[colback=red!5, colframe=red]
\textbf{Note}: 本表は \textbf{v8 として測定可能な最初の検証結果}.\\
v8.0 で得られた Phase 4 数値は Multi-framework 未機能による invalid な
測定であり、v8 の数値として扱わない.
\end{tcolorbox}

\begin{longtable}{l r r r r}
\toprule
Cell & v7.1 median & v8.1 median & paired diff & 判定 \\
\midrule
Normal\_H200 & ? & ? & ? & ? \\
Normal\_H500 & ? & ? & ? & ? \\
\textbf{Vulnerable\_H200} & 1.422 & 1.546 & \textcolor{success}{\textbf{+0.124}} & ✓ \\
\textbf{Vulnerable\_H500} & 1.347 & 1.546 & \textcolor{success}{\textbf{+0.199}} & ✓ \\
\textbf{Stagnation\_H200} & 14.146 & 3.140 & \textcolor{warning}{\textbf{-11.006}} & ⚠ 大悪化 \\
\textbf{Stagnation\_H500} & 13.806 & 3.140 & \textcolor{warning}{\textbf{-10.666}} & ⚠ 大悪化 \\
\bottomrule
\end{longtable}

\begin{honest}
\textbf{Multi-framework が機能した結果としての副作用}:

\begin{itemize}[leftmargin=*]
  \item Vulnerable で改善 (+0.124, +0.199) — NRMO 本旨が機能
  \item \textbf{Stagnation で大悪化 (-11)} — v8.1 が \texttt{invest/C} を選び早期破滅
  \item v8.1 は v7.1 を \textbf{安定して上回る性能には達していない}
\end{itemize}

つまり v8.1 は v7.1 に対し、世界によって改善と悪化が分かれる.
これは Multi-framework が \textbf{実際に意思決定に影響している証拠}.

ただし、世界別の制御が必要. 例えば Stagnation では Multi-framework の意見をそのまま採用せず、
World belief で重み付けする等の追加対応が必要.
\end{honest}

\section{V8.1 Long Run (paired)}

\begin{longtable}{l r r r}
\toprule
World & v7.1 time\_to\_ruin & v8.1 time\_to\_ruin & diff \\
\midrule
Normal & 52 & 13 & \textcolor{warning}{-39} \\
Vulnerable & 5 & 7 & \textcolor{success}{+2} \\
\bottomrule
\end{longtable}

\textbf{Vulnerable で +2 step 長く生存} (P1-2 の新基準で C3 PASS).

\section{V8.1 Phase 6 Final Judgment}

\begin{summary}
\textbf{判定}: PARTIAL\_PASS (4/5 PASS)

\begin{itemize}[leftmargin=*]
  \item \textcolor{warning}{✗ C1 PARETO\_IMPROVEMENT (2/6 = 33\%、閾値 90\% 未達)}
  \item \textcolor{success}{✓ C2 STRICT\_IMPROVEMENT (2 cells で改善)}
  \item \textcolor{success}{✓ C3 LONG\_RUN\_SAFETY (Vulnerable で実質改善)}
  \item \textcolor{success}{✓ C4 BLUE\_OCEAN (14 新規価値次元)}
  \item \textcolor{success}{✓ C5 V8\_INTEGRATION (14 レイヤー pipeline)}
\end{itemize}
\end{summary}

\section{v8.0 の Phase 4 数値の取扱い (重要な自己訂正)}

\begin{critical}
\textbf{重要な訂正}: v8.0 の Phase 4 数値は \textbf{v8 として無効} である.

理由:\\
v8.0 は Multi-framework key mismatch により、Multi-framework が
\textbf{実際には機能していなかった}.\\
従って v8.0 として記録された Phase 4 数値は、実態としては
\textbf{「legacy V71Engine + 一部 layer の影響」}であり、
\textbf{v8 アーキテクチャの性能を反映していない}.

そのような数値を「v8.0 の結果」として v8.1 と比較表に並べることは、
\textbf{測定対象が異なる数値を同じ軸で比較する数値不正}である.

\textbf{訂正}:\\
v8.0 Phase 4 数値: \textbf{archived / invalid as v8 measurement}\\
v8 の真の検証: v8.1 が \textbf{最初の有効な測定}
\end{critical}

\section{正しい比較: v7.1 (legacy) vs v8.1 (真の v8)}

\begin{longtable}{l r r r r}
\toprule
Cell & v7.1 median & v8.1 median & paired diff & 判定 \\
\midrule
Normal\_H200 & ? & ? & ? & ? \\
Normal\_H500 & ? & ? & ? & ? \\
\textbf{Vulnerable\_H200} & 1.422 & 1.546 & \textcolor{success}{\textbf{+0.124}} & ✓ \\
\textbf{Vulnerable\_H500} & 1.347 & 1.546 & \textcolor{success}{\textbf{+0.199}} & ✓ \\
\textbf{Stagnation\_H200} & 14.146 & 3.140 & \textcolor{warning}{\textbf{-11.006}} & ⚠ 大悪化 \\
\textbf{Stagnation\_H500} & 13.806 & 3.140 & \textcolor{warning}{\textbf{-10.666}} & ⚠ 大悪化 \\
\bottomrule
\end{longtable}

これが \textbf{v8 として測定可能な最初の数値} である.\\
Pareto pass rate 33\% (2/6) は v8.1 の真実の値.

\begin{honest}
\textbf{v8 の数値として}:

\begin{itemize}[leftmargin=*]
  \item Vulnerable で改善 (+0.124, +0.199) — NRMO 本旨が機能
  \item \textbf{Stagnation で大悪化 (-11)} — Multi-framework が \texttt{invest/C} を選び早期破滅
\end{itemize}

これは v8.1 が v7.1 に対し:
\begin{itemize}[leftmargin=*]
  \item 一部の世界では改善
  \item 一部の世界では大きく悪化
\end{itemize}
を示す. v8 architecture は \textbf{まだ v7.1 を安定して上回るレベルにない}.
\end{honest}

% ============================================================
\chapter{Stagnation 対策の試行と撤回 (v8.2 実験記録)}

\section{動機}

v8.1 の paired Phase 4 で Stagnation\_H200/H500 において
\textbf{-11.006 / -10.666 の大悪化} が観測された.\\
原因: v8.1 の Multi-framework が \texttt{invest/C} を選び、Stagnation で resource を浪費し早期破滅.

\textbf{ユーザー指示}: 「スタグネーション対策を」

\section{試行 1: Belief-based 検出}

\textbf{設計}: \texttt{BayesianUpdater.belief.mean()} から
\texttt{opportunity\_arrival\_rate} と \texttt{opportunity\_window\_duration} を取得し、
両方低ければ Stagnation と判定.

\textbf{結果}: \textbf{失敗}.\\
50 step 観測しても belief mean が真値 (rate=0.05, window=20) に収束せず、
mean estimate は rate≈0.19, window≈10. Particle Filter の精度不足.\\
\texttt{stagnation\_score = 0.0} で trigger せず.

\section{試行 2: 観測ベース (state.O + 平均 reward)}

\textbf{設計}: state.O が低い + 直近 reward 平均が低い → Stagnation.

\textbf{結果}: \textbf{失敗}.\\
Stagnation world でも state.O は 47-52 で安定 (低くない).
平均 reward も 0.30 程度で大差なし.\\
私の事前仮定 (Stagnation = O が低い世界) が \textbf{完全な誤解} だった.\\
正しくは: Stagnation = 機会到来が稀 (rate=0.05) だが、来た機会の窓は長い (window=20).
状態 O は安定するが、resource decay が支配的.

\section{試行 3: invest 連発 + R 下降 (v8.2 ver.1)}

\textbf{設計}: 直近 10 step で invest > 50\% かつ R 下降 → invest を recover に切替.

\textbf{結果}: \textbf{部分成功} + \textbf{副作用}.\\

\begin{tabular}{l r r}
\toprule
Cell & v8.1 diff & v8.2 ver.1 diff \\
\midrule
Stagnation\_H200 (single trial) & -11 → \textbf{-9.7} & 改善 \\
Cumulative score on seed=42 & 3.30 → \textbf{7.81} & 改善 \\
Normal\_H500 (paired n=100) & ? & \textcolor{warning}{\textbf{-11.49}} \\
Pareto pass rate & 2/6 (33\%) & 1/6 (16.7\%) \\
\bottomrule
\end{tabular}

\textbf{副作用の原因}: Normal world でも v8.1 が invest 連発する.
Safeguard が誤発火し、invest/C を recover/A に切替 → score 大悪化.

\section{試行 4: 厳格化 (early phase 除外 + 2 条件以上要件) (v8.2 ver.2)}

\textbf{設計}: 早期 phase (≤8 step) は trigger しない + 2 条件以上同時必要.

\textbf{結果}: \textbf{さらに悪化}.\\
Pareto pass: 1/6 (16.7\%) \(\rightarrow\) \textbf{0/6 (0\%)}.\\
全 cell で v8.1 より悪化. 厳格化により Stagnation でも trigger せず、
かつ Normal でも trigger しない \texttt{recent\_actions} 履歴管理の副作用で他の動作も変化.

\section{結論: safeguard 撤回}

\begin{honest}
\textbf{表面対症療法では Multi-framework の根本バイアスを解決できない}.

v8.1 で Multi-framework が \texttt{invest/C} を選びすぎるという根本問題に対し、
事後的に safeguard で抑制しようとした.\\
しかし:

\begin{itemize}[leftmargin=*]
  \item belief 精度不足で世界を識別できない
  \item 観測ベースの heuristic は誤発火しやすい
  \item Stagnation だけ trigger する厳格化を作ろうとすると、結局複雑化して逆効果
\end{itemize}

\textbf{真の対策}: Multi-framework 自体の調整 (EUT 重み、世界推定との動的連動).

v8.2 の safeguard を \textbf{撤回} し、v8.1 アーキテクチャを維持.
v8.3 で根本対策を試みる.
\end{honest}

\section{v8.2 で行った付随変更}

safeguard 自体は撤回したが、以下は v8.2 として保持:

\begin{itemize}[leftmargin=*]
  \item \texttt{recent\_rewards}, \texttt{recent\_actions}, \texttt{recent\_R\_values} の
        履歴管理機構を engine に追加 (v8.3 で活用予定)
  \item \texttt{update\_reward()} 内での履歴 update
  \item 14/14 test PASS 維持
  \item Decision pipeline 14 レイヤー構造は v8.1 と同じ
\end{itemize}

\section{v8.2 の honest な数値}

\textbf{paired Phase 4 (safeguard 撤回後)}:

\begin{tabular}{l r r r}
\toprule
Cell & v7.1 median & v8.2 median & diff \\
\midrule
Stagnation\_H200 & 14.183 & 3.140 & \textcolor{warning}{-11.043} \\
Stagnation\_H500 & 13.994 & 3.140 & \textcolor{warning}{-10.854} \\
Pareto pass rate & — & 1/6 (16.7\%) & — \\
\bottomrule
\end{tabular}

これは v8.1 とほぼ同じ性能 (Stagnation 副作用は依然存在).\\
\textbf{v8.2 = Stagnation 対策の \emph{失敗の記録}.}\\
無理に「対策できた」と書かない. honest に「未解決」と記録する.

\section{v8.3 への学び}

\begin{summary}
\textbf{得られた洞察}:

\begin{enumerate}[leftmargin=*]
  \item Stagnation の本質は \textbf{機会到来稀 + 機会窓長い} (window 長は誤認していた)
  \item belief Particle Filter の精度は低い (50 step では収束しない)
  \item 観測ベース heuristic は誤発火リスクが高い
  \item 事後 safeguard でなく、Multi-framework 段階での調整が必要
\end{enumerate}

\textbf{v8.3 候補設計}:

\begin{itemize}[leftmargin=*]
  \item EUT の評価重みを動的に下げる (世界 belief が unclear な時)
  \item Multi-framework に Phase 9 Hyperbolic discounter を統合し、
        遠い期待値の影響を弱める
  \item Phase 9 Causal graph で「invest が R を減らす」因果効果を反映
  \item 履歴の R/E/O trajectories を Multi-framework input に追加
\end{itemize}
\end{summary}

% ============================================================
\chapter{残された課題と v8.3 への展望}

\section{v8.1 で未対応の項目}

\subsection{P2-1: placeholder layer 実効化}

現在 placeholder のまま:
\begin{itemize}[leftmargin=*]
  \item Goodhart metrics は実際の reward/optionality から計算していない
  \item Meta-cognition は予測と実 outcome の照合が次 step で行われていない
  \item External feedback は記録のみ、次回 decision に影響していない
  \item Tower distance は出力されているが、confidence に反映されていない
\end{itemize}

\subsection{Stagnation 副作用への対応}

Multi-framework が \texttt{invest/C} を選ぶ場面で、World belief が
「Stagnation である」と認識した場合に override する仕組みが必要.

\subsection{World simulation 自体の見直し}

依然として両エンジンで ruin\_rate 100\%. これは v8 の問題ではなく World simulation の
自然動態. ruin 判定の再設計または短期 horizon 化が必要.

\subsection{Phase 7-9 機構の更なる V8Engine 活用}

\begin{itemize}[leftmargin=*]
  \item CMA-ES で候補ランキング
  \item Causal graph で counterfactual 予測
  \item Meta-cog で confidence 動的調整
  \item Survivorship correction で action value 補正
\end{itemize}

\section{v8.2 への候補ロードマップ}

\begin{enumerate}[leftmargin=*]
  \item World-conditional Multi-framework weighting\\
        Belief 上の World 推定に応じて framework 重みを変える
  \item Goodhart metric の実計算化\\
        Score の trend と他指標の trend の divergence で発火
  \item Meta-cog 動作\\
        前 step の予測と実 outcome を照合、confidence を補正
  \item External feedback の動作化\\
        ユーザーフィードバックを reward/threshold に反映
  \item Tower distance → confidence reduction
  \item Anti-fragile barbell の実 layer 化\\
        80\% safe + 20\% risky の resource 配分
  \item Adversarial training\\
        worst-case world での訓練
\end{enumerate}

% ============================================================
\chapter{結語}

\section{監査と honest な評価}

第三者監査により v8.0 の最大の隠れバグ (Multi-framework key mismatch) が露呈し、
v8.1 で修正された.

\begin{summary}
\textbf{v8.1 の honest な現状}:

\begin{itemize}[leftmargin=*]
  \item ✓ 構造統合は \textbf{真に} 動作している
  \item ✓ 再現性は assertion で保証
  \item ✓ 14/14 test PASS
  \item ✗ Pareto pass 2/6 (33\%) — v7.1 を安定して上回らない
  \item ✗ Stagnation で大きな副作用
\end{itemize}

\textbf{完成度評価 (honest)}:\\
v8.0 と v8.1 を「完成度」で比較すること自体が誤り.\\
v8.0 は v8 アーキテクチャを反映していなかった (Multi-framework 未機能).\\
v8.1 は v8 として\textbf{初めて測定可能}になった.\\
従って完成度は v8.1 を基準に: \textbf{Architecture candidate stage}.
\end{summary}

\section{次の世代へ}

v8 architecture の \textbf{方向性は正しい}.\\
ただし、より honest であるためには、まだ調整が必要.

\begin{tcolorbox}[colback=navy!5, colframe=navy, width=\linewidth]
\centering
\textbf{呼称}:\\
v8.0: v8 部品を持つ v7.2 (前回監査の判定)\\
v8.1: v8 runtime prototype (今回)\\
v8.2 候補: World-conditional control + placeholder 実効化\\
v8 release candidate: Pareto pass rate 90\% 達成後
\end{tcolorbox}

\vspace{1cm}

\rule{\linewidth}{0.5pt}

\centering
\emph{NRMO v8.1 Patch Report — 完}

\end{document}
